\documentclass[10pt,a4paper]{article}
\usepackage{amsmath,amssymb,amsthm,graphicx}
\usepackage{amsmath,amssymb,amsthm}
\newcommand{\supp}{\operatorname{supp}}
\newcommand{\Newt}{\operatorname{Newt}}
\newcommand{\conv}{\operatorname{conv}}
\newcommand{\vertset}{\operatorname{vert}}
\newenvironment{sourceresult}[1]{\par\medskip\noindent\textbf{#1.}\itshape}{\par\medskip}
\title{Linear-size monotone arithmetic circuits with exponentially many Newton polygon vertices}
\author{Pavel Hrubeš, Amir Yehudayoff}
\date{doi:10.4230/LIPIcs.CCC.2021.9; CC BY 4.0}
\begin{document}
\maketitle
\noindent\textit{Editorial note.} Faithful editable transcription of the official CC BY publication; full construction and bearing author arguments retained with original numbers and page locators. Original displayed algebra and wording retained without mathematical correction. Long display lines mechanically split for page width. Journal alias identifies the same work but no journal-only content is used. Production and validation solely on Yupeng.\par\medskip
\noindent\textit{Allowed standard prerequisites (editorial).}\begin{itemize}
\item Convex hulls, finite convexly independent sets and arithmetic circuits. Standard background; original shadow, transparency, h.p.-circuit and h.p.-projection definitions retained. Folklore planar Minkowski-sum lemma retained with author outline and original cited source.
\end{itemize}\medskip
\section*{Original definitions and computational models (pp.9:1--2,9:11--14)}
A polytope is the convex hull of a finite set of points in Euclidean space. Equivalently, it is a compact set that is defined by finitely many linear inequalities.
A shadow of a polytope $P\subseteq\mathbb R^n$ is a set of the form $L(P)$, where $L:\mathbb R^n\to\mathbb R^2$ is a linear map. The shadows of $P$ are two-dimensional polygons, and hence typically much simpler than $P$. The shadow complexity of $P$ is
\[
\sigma(P)=\max_L |\vertset(L(P))|,
\]
where $L$ is a linear map and $\vertset(Q)$ is the vertex set of the polytope $Q$.
Let $\mathbb F$ be a field. For a list of variables $x=(x_1,\ldots,x_n)$ and $\alpha\in\mathbb N^n$, let $x^\alpha$ be the monomial $\prod_{i=1}^n x_i^{\alpha_i}$. A polynomial $f\in\mathbb F[x_1,\ldots,x_n]$ is a formal sum of the form $\sum_{\alpha\in\mathbb N^n}a_\alpha x^\alpha$ where $\supp(f):=\{\alpha\in\mathbb N^n:a_\alpha\ne0\}$ is finite. The Newton polytope of $f$ is
\[\Newt(f):=\conv(\supp(f)),\]
where $\conv(\cdot)$ denotes the convex hull.
A polynomial $f$ is transparent if $|\supp(f)|=\sigma(\Newt(f))$. In other words, there is a linear map $L:\mathbb R^n\to\mathbb R^2$ which maps $\supp(f)$ to distinct convexly independent points in $\mathbb R^2$.

As the standard model of computation of polynomials we take the arithmetic circuit model. It is a (finite) directed acyclic graph whose every node has in-degree zero or two. Nodes of in-degree zero (input nodes) are labelled with a constant or a variable. Nodes of in-degree two are labelled with operations $+$ or $\times$. Every node of the circuit computes a polynomial in $\mathbb F$ in the natural way. As the size of the circuit, we take the number of nodes. A circuit is called a formula if its underlying graph is a tree. For more background and motivation, see [44] and references within.
Our focus is mainly on monotone computation. A polynomial over $\mathbb R$ is monotone if all of its coefficients are non-negative. Similarly, a monotone arithmetic circuit can use only non-negative constants.

In a Laurent polynomial, variables are allowed to have negative integer exponents. The notions of $\supp(f)$ and Newton polytope of $f$ are defined in the obvious manner. A Laurent polynomial over $\mathbb R$ is monotone, if all of its coefficients are non-negative.
A high powered circuit (h.p.-circuit for short) is an arithmetic circuit in which every input node is labelled by a term $ax_1^{k_1}\cdots x_n^{k_n}$ with $a\in\mathbb F$ and $k_1,\ldots,k_n\in\mathbb Z$. The size of the circuit is the number of its nodes.
In other words, we have given the circuit a power to compute every term $ax^\alpha$ at a unit cost. This is especially important in the case of h.p.-formula. An arithmetic formula of size $s$ can compute a polynomial of degree at most $s$, whereas there is no such restriction in an h.p.-formula. Furthermore, we have allowed the variables to have negative exponents and hence an h.p.-circuit computes a Laurent polynomial instead of a polynomial. But this is only a cosmetic detail.

A high power projection (h.p.-projection for short) is a homomorphism
\[\pi:\mathbb F[x_1,\ldots,x_n]\to\mathbb F[y_1,\ldots,y_m,y_1^{-1},\ldots,y_m^{-1}]\]
such that $\pi(x_i)=a_i y^{\alpha_i}$ for every $x_i$, where $a_i\in\mathbb F$ and $\alpha_i\in\mathbb Z^m$, and for every $f\in\mathbb F[x_1,\ldots,x_n]$,
\[\pi(f(x_1,\ldots,x_n))=f(\pi(x_1),\ldots,\pi(x_n)).\]
The constants $a_i$ are called the coefficients of $\pi$ and $\alpha_i$ its exponents. If $\mathbb F=\mathbb R$ and $\pi$ has non-negative coefficients, then $\pi$ is called monotone.
An h.p.-projection $\pi$ induces a linear map $L_\pi:\mathbb R^n\to\mathbb R^m$ by setting $L_\pi(e_i)=\alpha_i$ and extending it linearly to $\mathbb R^n$. It follows that $\supp(\pi(f))\subseteq L_\pi(\supp(f))$. The inclusion may be strict, as some monomials of $f$ can cancel out in the projection. If no cancellations occur, we indeed have $\Newt(\pi(f))=L_\pi(\Newt(f))$. This is satisfied, e.g., if $f$ is monotone and the coefficients of $\pi$ are positive, or if the coefficients are algebraically independent.

\begin{sourceresult}{Theorem 30 (p.9:13)}
For every $n$, there exists a monotone bivariate polynomial $f_n$ such that $f_n$ has a monotone arithmetic circuit of size $O(n)$ and $\Newt(f)$ has $2^n$ vertices.
\end{sourceresult}

\section*{Original bearing operations (pp.9:6--7)}
Given $A,B\subseteq\mathbb R^n$, their Minkowski sum is defined as
\[A+B:=\{a+b:a\in A,b\in B\}.\]
If $P$ and $Q$ are polytopes then $P+Q$ is also a polytope. In two-dimensions, Minkowski sum has nice properties. Let $P$ be a polytope in $\mathbb R^2$ with vertices $v_1,\ldots,v_k$ where $k>1$. We can assume they are ordered so that $P$ lies in the left closed half plane determined by the line going from $v_i$ to $v_{i+1}$ for $i<k$, and similarly for $v_k$ and $v_1$. Let $E(P)$ be the collection of unit vectors in the direction of these $k$ edges. That is, vectors of the form $(v_{i+1}-v_i)/\|v_{i+1}-v_i\|$ for $i<k$, and $(v_1-v_k)/\|v_1-v_k\|$. If $|\vertset(P)|\le1$ then $E(P):=\emptyset$.
\begin{sourceresult}{Lemma 8}
Let $P_1,\ldots,P_r$ be non-empty polytopes in $\mathbb R^2$. Then $E(P_1+\cdots+P_r)=\bigcup_{i=1}^r E(P_i)$. Consequently, $|\vertset(P_1+\cdots+P_r)|\le\sum_{i=1}^r|\vertset(P_i)|$. The latter holds for empty $P_i$'s as well.
\end{sourceresult}
The lemma is folklore. It can be inferred from Chapter 13.3 in [11], and we give only an outline of proof.
\begin{proof}[Proof sketch of Lemma 8]
Given a non-empty polytope $P$ and $w\in\mathbb R^2$, let $P^w:=\{x\in P:\langle x,w\rangle=\max_{z\in P}\langle z,w\rangle\}$ be the set of extreme of points of $P$ in the direction $w$. It is either a vertex or an edge of $P$. For a pair of polytopes we have $(P_1+P_2)^w=P_1^w+P_2^w$. Every edge of $P_1$ yields an edge of $P_1+P_2$ with the same direction. Conversely, every edge of $P_1+P_2$ comes from one of $P_1$ or $P_2$.
\end{proof}
The second operation we use is
\[A\sqcup B:=\conv(A\cup B).\]
If $P$ and $Q$ are polytopes then $P\sqcup Q$ is also a polytope.
\begin{sourceresult}{Lemma 11}
Let $f,g$ be Laurent polynomials over $\mathbb F$.
\begin{enumerate}
\item[(i)] Then $\Newt(fg)=\Newt(f)+\Newt(g)$.
\item[(ii)] $\Newt(f+g)=\Newt(f)\sqcup\Newt(g)$, provided $\mathbb F=\mathbb R$ and both $f$ and $g$ are monotone.
\end{enumerate}
\end{sourceresult}
\begin{proof}
Part (i) can be found in [16] for polynomials; it extends to Laurent polynomials. Part (ii) is straightforward to verify.
\end{proof}

\section*{Entire original Section 8: An easy polynomial with many vertices (pp.9:19--21)}
Here we construct a bivariate polynomial with a monotone arithmetic circuit of linear size, but whose Newton polytope has exponentially many vertices. This proves Theorem 30.
We use the following notation. Given $(a,b)\in\mathbb R^2$,
\[(a,b)\cdot P:=\{(ax,by):(x,y)\in P\}.\]
Given $a\in\mathbb R$,
\[aP:=(a,a)\cdot P.\]
\begin{sourceresult}{Observation 52}
For a bivariate polynomial $f(x,y)$,
\[\Newt(f(x^a,y^b))=(a,b)\Newt(f(x,y))\quad\text{and}\quad\Newt(f^a)=a\Newt(f).\]
\end{sourceresult}
The building block of the polynomial are the following two polytopes. Let $P_n$ be the polytope with vertices $\{(k,k^2):0\le k\le n-1\}$. Let $Q_n$ be the polytope with vertices $\{(k,k^2+k):0\le k\le n-1\}$. These polytopes can be constructed inductively as follows.
\begin{sourceresult}{Lemma 53}
For every $n\ge1$,
\begin{align}
P_{2n}&=(2,4)\cdot P_n\sqcup((1,1)+(2,4)\cdot Q_n)),\tag{8}\\
Q_{2n}&=(1,2)\cdot(P_n+Q_n)\sqcup\{(2n-1,2n(2n-1))\}.\tag{9}
\end{align}
\end{sourceresult}
\begin{proof}
Part (8). Let $0\le k\le2n-1$. If $k=2r$ is even then $r\le n-1$ and
\[(k,k^2)=(2,4)(r,r^2)\]
with $(r,r^2)$ a vertex of $P_n$. If $k=2r+1$ is odd then $r\le n-1$ and
\[(k,k^2)=(2r+1,4r^2+4r+1)=(1,1)+(2,4)\cdot(r,r^2+r),\]
where $(r,r^2+r)$ is a vertex of $Q_n$. This shows the containment $\subseteq$ in (8). The other direction holds since $P_n\sqcup Q_n$ can have at most $2n$ vertices.

Part (9). We first describe the vertices of $(1,2)(P_n+Q_n)$. We claim that
\begin{equation}
\vertset((1,2)(P_n+Q_n))=\{v_0,v_1,\ldots,v_{2n-2},u\},\tag{10}
\end{equation}
where $v_k:=(k,k^2+k)$, $u:=(n-1,2n(n-1))$.
Given $0\le k\le2n-2$, let us show that $v_k$ is a vertex of $(1,2)(P_n+Q_n)$. If $k=2r$ is even, we have $r\le n-1$ and
\[(k,k^2+k)=(2r,4r^2+2r)=(1,2)(r,r^2)+(1,2)(r,r^2+r).\]
If $k=2r+1$ is odd, we have $r\le n-2$ and
\[(k,k^2+k)=(2r+1,4r^2+6r+2)=(1,2)(r+1,(r+1)^2)+(1,2)(r,r^2+r).\]
This means that $v_k\in(1,2)(P_n+Q_n)$. Now, every $(z_1,z_2)\in(1,2)(P_n+Q_n)$ satisfies $z_2\ge z_1^2+z_1$, because
\[2r_1^2+2(r_2^2+r_2)-(r_1+r_2)^2-(r_1+r_2)=(r_1-r_2)^2-(r_1-r_2)\ge0.\]
Since $v_k$ lies on the curve $z_2=z_1^2+z_1$, and the curve is strictly convex, $v_k$ cannot be convex combination of other points in $(1,2)(P_n+Q_n)$. So, $v_k$ is indeed a vertex. To show that $u$ is a vertex, note that both $(1,2)P_n$ and $(1,2)Q_n$ are contained in the halfplane $\{(z_1,z_2)\in\mathbb R^2:z_2\le2nz_1\}$. On the boundary $z_2=2nz_1$, $(1,2)Q_n$ has vertices $(0,0)$ and $u$, and $(1,2)P_n$ only the vertex $(0,0)$. This implies $u$ is a vertex of $(1,2)(P_n+Q_n)$. This proves the containment $\subseteq$ in (10). Equality holds since $P_n+Q_n$ can have at most $2n$ vertices.
To infer (9) from (10), note that $u$ lies on the line connecting the origin and $v_{2n-1}=(2n-1,2n(2n-1))$.
\end{proof}
\begin{proof}[Proof of Theorem 30]
Inductively define a sequence of bivariate polynomials. The base case is
\[p_0=1\quad\text{and}\quad q_0=1.\]
The inductive step is
\[p_{n+1}=p_n(x^2,y^4)^2+x^Ny^Nq_n(x^2,y^4)^2\]
and
\[q_{n+1}=p_n(x^2,y^4)q_n(x^2,y^4)+x^{N(N-1)}y^{N^2(N-1)}\]
where $N=2^{n+1}$.
We claim that for every $n\ge0$,
\begin{equation}
\Newt(p_n)=2^nP_{2^n}\quad\text{and}\quad\Newt(q_n)=2^nQ_{2^n}.\tag{11}
\end{equation}
For $n=0$, this follows from $\Newt(p_0)=\Newt(q_0)=\{(0,0)\}=P_1=Q_1$. The induction step uses Lemma 11 and Observation 52. Assume that (11) holds for a given $n\ge0$. Then
\[\Newt(p_n(x^2,y^4))=2^n(2,4)P_{2^n}\quad\text{and}\quad\Newt(q_n(x^2,y^4))=2^n(2,4)Q_{2^n}.\]
Using (8),
\begin{align*}
\Newt(p_{n+1})&=2\cdot2^n(2,4)P_{2^n}\sqcup((N,N)+2\cdot2^n(2,4)Q_{2^n}))\\
&=2^{n+1}((2,4)P_{2^n}\sqcup((1,1)+(2,4)Q_{2^n})))\\
&=2^{n+1}P_{2^{n+1}}.
\end{align*}
Similarly, part (9) gives
\begin{align*}
\Newt(q_{n+1})&=2^n(2,4)(P_{2^n}+Q_{2^n})\sqcup\{(N(N-1),N^2(N-1))\}\\
&=2^{n+1}\big((1,2)(P_{2^n}+Q_{2^n})\sqcup\{(N-1,N(N-1))\}\big)\\
&=2^{n+1}Q_{2^{n+1}}.
\end{align*}
This proves (11).
To compute $p_n,q_n$, first construct a circuit of size $O(n)$ that simultaneously computes $x^M,x^{M(M-1)},y^M,y^{M^2(M-1)}$ for every $M=2^m$ with $m\le n$. Now, construct a circuit for $p_n$ and $q_n$ inductively. Given a circuit for $p_n$ and $q_n$, we can construct a new one computing $p_{n+1},q_{n+1}$ by introducing a constant number of extra gates.
\end{proof}

\begin{thebibliography}{44}
\bibitem[11]{source11} M. de Berg, M. van Kreveld, M. Overmars, and O. Schwarzkopf. \emph{Computational Geometry: Algorithms and Applications}. Springer, 2 edition, 2000.
\bibitem[44]{source44} A. Shpilka and A. Yehudayoff. Arithmetic circuits: A survey of recent results and open questions. \emph{Foundations and Trends in Theoretical Computer Science}, 5(3-4), 2010.
\bibitem[16]{source16} S. Gao. Absolute irreducibility of polynomials via newton polytopes. \emph{Journal of Algebra}, 237(2):501--520, 2001.
\end{thebibliography}
\end{document}
